\documentclass[11pt]{article}

\usepackage[margin=1in]{geometry}
\usepackage{amsmath,amssymb,amsthm}
\usepackage{graphicx}
\usepackage{hyperref}
\usepackage{microtype}

\newtheorem{theorem}{Theorem}
\newtheorem{remark}{Remark}
\newcommand{\Fix}{\operatorname{Fix}}
\newcommand{\Lip}{\operatorname{Lip}}

\title{A Hilbert-Space Partial Answer to a Uniformly Lipschitz Fixed-Point Question}
\author{Smoke-test packet for arXiv:2307.12958}
\date{June 11, 2026}

\begin{document}
\maketitle

\section*{Source Paper}

This packet concerns the paper
\[
\textit{Retraction methods and fixed point free maps with null minimal
displacements on unit balls}
\]
by Cleon S. Barroso and Valdir Ferreira, arXiv:2307.12958.

The arXiv source paper asks the following introductory questions after
recalling Lin--Sternfeld's positive minimal displacement theorem.

\begin{center}
\includegraphics[width=\textwidth]{figures/open_problem_crop.png}
\end{center}

The review thread also supplied the following screenshot of a
journal/published formulation.  A separate publicly indexed journal PDF was
not located; this packet therefore keeps the arXiv PDF as the source paper and
includes this screenshot to document the alternate wording.

\begin{center}
\includegraphics[width=\textwidth]{figures/journal_question_crop.png}
\end{center}

\section*{Target}

We focus on the second question.  In the arXiv text it is stated as:
\begin{quote}
\((Q2)\) Is there a uniformly Lipschitz map \(T:K\to K\) with
\(\Fix(T)=\varnothing\)?
\end{quote}

The journal screenshot asks the stronger unit-ball version:
\begin{quote}
Is there a uniformly Lipschitz mapping \(T:B_X\to B_X\) satisfying
\(\mathrm d(T,B_X)=0\) and \(\Fix(T)=\varnothing\)?
\end{quote}

The result below answers this stronger version for infinite-dimensional
Hilbert spaces.  It is a partial result only; it does not settle the general
Banach-space question.

\section*{Candidate Partial Answer}

\begin{theorem}
Let \(H\) be an infinite-dimensional real Hilbert space.  Then there exists
a map \(T:B_H\to B_H\) such that
\[
\Fix(T)=\varnothing,\qquad
\mathrm d(T,B_H)=0,\qquad
\sup_{n\ge 1}\Lip(T^n)<\infty .
\]
\end{theorem}

\begin{proof}
Choose a unit vector \(e_0\in H\).  Since \(H\) is infinite-dimensional,
there is a linear isometry
\[
U:H\to e_0^\perp
\]
with no nonzero fixed vector.  In the separable model this is the unilateral
shift \(Ue_j=e_{j+1}\).  In the nonseparable case one may use such a shift on
a separable summand and \(-I\) on the remaining orthogonal summand.

Fix a Lipschitz function \(\theta:[0,1]\to[0,\pi/2]\) such that
\[
\theta(r)=0\quad(0\le r\le 1/2),\qquad
\theta(1)=\pi/2 .
\]
For instance, take \(\theta(r)=0\) on \([0,1/2]\) and
\(\theta(r)=\pi(r-1/2)\) on \([1/2,1]\).  Define
\[
a(r)=\cos\theta(r),\qquad
b(r)=
\begin{cases}
\dfrac{\sin\theta(r)}{r},& r>0,\\[2mm]
0,& r=0.
\end{cases}
\]
Because \(\theta\) vanishes on a neighborhood of \(0\), both \(a\) and \(b\)
are Lipschitz on \([0,1]\).

For \(x\in B_H\), set \(r=\|x\|\) and define
\[
T(x)=a(r)e_0+b(r)Ux .
\]
Since \(e_0\perp Ux\), we have
\[
\|T(x)\|^2=a(r)^2+b(r)^2\|x\|^2
=\cos^2\theta(r)+\sin^2\theta(r)=1 .
\]
Thus \(T\) maps \(B_H\) into the unit sphere \(S_H\).

The map \(T\) is Lipschitz.  Indeed, \(x\mapsto \|x\|\) is
1-Lipschitz, \(a\) and \(b\) are Lipschitz, \(b\) is bounded, and \(U\) is an
isometry.  Hence, for \(x,y\in B_H\),
\[
\begin{aligned}
\|T(x)-T(y)\|
&\le |a(\|x\|)-a(\|y\|)|\\
&\quad+\|b(\|x\|)Ux-b(\|y\|)Uy\|\\
&\le \bigl(\Lip(a)+\|b\|_\infty+\Lip(b)\bigr)\|x-y\|.
\end{aligned}
\]

On the unit sphere, \(\theta(1)=\pi/2\), so \(a(1)=0\) and \(b(1)=1\).
Therefore
\[
Tz=Uz\qquad(z\in S_H).
\]
Since \(T(B_H)\subset S_H\), an induction gives
\[
T^n x=U^{\,n-1}T x\qquad(x\in B_H,\ n\ge 1).
\]
Because \(U\) is an isometry,
\[
\Lip(T^n)\le \Lip(T)\qquad(n\ge 1).
\]
Thus \(T\) is uniformly Lipschitz.

There is no fixed point.  If \(T(x)=x\), then \(x\in S_H\), so
\[
x=T(x)=Ux .
\]
This contradicts the choice of \(U\), which has no nonzero fixed vector.

It remains to check null minimal displacement.  In the separable shift
summand, let
\[
v_N=\frac{e_0+e_1+\cdots+e_{N-1}}{\sqrt N}\in S_H .
\]
Then \(T(v_N)=Uv_N\), and hence
\[
\|T(v_N)-v_N\|
=\left\|\frac{e_N-e_0}{\sqrt N}\right\|
=\sqrt{\frac{2}{N}}\longrightarrow 0 .
\]
Therefore \(\mathrm d(T,B_H)=0\).
\end{proof}

\section*{Why This Does Not Settle the General Q2}

The proof uses a Hilbert-space isometry \(U:H\to e_0^\perp\) with no fixed
vector and with approximate fixed points.  This shift geometry is exactly the
kind of structure that is missing in an arbitrary Banach space.  The packet
therefore gives a genuine special-case solution to the stronger journal-style
Q2, but it does not remove the paper's general obstruction.

\begin{remark}
A tempting older Hilbert map is
\[
x\mapsto (1-\|x\|,x_1,x_2,\ldots).
\]
It is Lipschitz and fixed-point free, but its iterates are not uniformly
Lipschitz: near the sphere the radial deficit decreases only quadratically,
so sensitivity can accumulate over about \(1/(1-\|x\|)\) iterates.  The collar
construction above avoids this by landing on the sphere in one step; after
that all iterates are isometric shifts.
\end{remark}

\section*{Human Review Recommendation}

Review as a partial answer to \((Q2)\), not as a full answer to the general
Banach-space problem.  The key checks are that \(b\) is Lipschitz at \(0\),
that \(T(B_H)\subset S_H\), and that \(T^n=U^{n-1}T\) for \(n\ge 1\).

\end{document}

